\subsection{Anillos ádicos y límites proyectivos}
\label{subsection:0.7.2-fr}

\begin{env}[7.2.1]
\label{0.7.2.1-fr}
Todo límite proyectivo de anillos \emph{discretos} es evidentemente un anillo
linealmente topologizado, separado y completo. Recíprocamente, sea $A$ un anillo
linealmente topologizado, y sea $(\mathfrak J_\lambda)$ un sistema fundamental
de vecindades abiertas de $0$ en $A$ formado

\oldpage[0\textsubscript{I}]{63}

por ideales. Las aplicaciones canónicas
$\varphi_\lambda:A\to A/\mathfrak J_\lambda$ forman un sistema proyectivo de
representaciones continuas y definen, por tanto, una representación continua
$\varphi:A\to\varprojlim A/\mathfrak J_\lambda$; si $A$ es \emph{separado},
$\varphi$ es un isomorfismo topológico de $A$ sobre un
subanillo denso en $\varprojlim A/\mathfrak J_\lambda$; si, además,
$A$ es \emph{completo}, $\varphi$ es un isomorfismo topológico de $A$ \emph{sobre}
$\varprojlim A/\mathfrak J_\lambda$.
\end{env}

\begin{lemma}[7.2.2]
\label{0.7.2.2-fr}
Para que un anillo linealmente topologizado sea admisible, es necesario y suficiente
que sea isomorfo a un límite proyectivo
$A=\varprojlim A_\lambda$, donde $(A_\lambda,u_{\lambda\mu})$ es un sistema
proyectivo de anillos discretos que tiene como conjunto de índices un conjunto ordenado
filtrante $L$ (para $\leq$) que admite un menor elemento denotado por $0$ y satisface
las condiciones siguientes: 1\textsuperscript{o} los
$u_\lambda:A\to A_\lambda$ son sobreyectivos; 2\textsuperscript{o} el núcleo
$\mathfrak J_\lambda$ de $u_{0\lambda}:A_\lambda\to A_0$ es nilpotente.
Cuando esto ocurre, el núcleo $\mathfrak J$ de $u_0:A\to A_0$ es igual a
$\varprojlim\mathfrak J_\lambda$.
\end{lemma}

La necesidad de la condición resulta de (7.2.1), tomando como
$(\mathfrak J_\lambda)$ un sistema fundamental de vecindades de $0$ formado
por ideales de definición contenidos en uno de ellos $\mathfrak J_0$ y aplicando
(7.1.4, (i)). El recíproco resulta de la definición de un límite proyectivo y
de (7.1.2), y la última afirmación es inmediata.

\begin{env}[7.2.3]
\label{0.7.2.3-fr}
Sean $A$ un anillo topológico \emph{admisible}, $\mathfrak J$ un ideal de $A$
contenido en un ideal de definición (es decir (7.1.4), tal que
$(\mathfrak J^n)$ tiende a $0$); se puede considerar sobre $A$ la topología
de anillo que tiene como sistema fundamental de vecindades de $0$ las potencias
$\mathfrak J^n$ ($n>0$); la llamaremos también la topología
\emph{$\mathfrak J$-preádica}. La hipótesis de que $A$ es admisible implica que
$\bigcap_n\mathfrak J^n=0$, por tanto la topología $\mathfrak J$-preádica sobre $A$
es \emph{separada}; sea
$\widehat A=\varprojlim A/\mathfrak J^n$ el completado de $A$ para esta topología
(donde los $A/\mathfrak J^n$ están provistos de la topología discreta), y designemos
por $u$ el homomorfismo de anillos $A\to\widehat A$ (no necesariamente
continuo), límite proyectivo de la sucesión de homomorfismos
$u_n:A\to A/\mathfrak J^n$. Por otra parte, la topología $\mathfrak J$-preádica
sobre $A$ es más fina que la topología dada $\mathcal T$ sobre $A$; como $A$
es separado y completo para $\mathcal T$, se puede prolongar por continuidad
la aplicación identidad de $A$ (provisto de la topología $\mathfrak J$-preádica)
en $A$ provisto de $\mathcal T$; esto da una representación continua
$v:\widehat A\to A$.
\end{env}

\begin{proposition}[7.2.4]
\label{0.7.2.4-fr}
Si $A$ es un anillo admisible y $\mathfrak J$ está contenido en un ideal de
definición de $A$, $A$ es separado y completo para la topología
$\mathfrak J$-preádica.
\end{proposition}

En efecto, con las notaciones de (7.2.3), es inmediato que $v\circ u$ es
la aplicación identidad de $A$. Por otra parte,
$u_n\circ v:\widehat A\to A/\mathfrak J^n$ es la prolongación por continuidad
(para la topología $\mathfrak J$-preádica sobre $A$ y la topología discreta sobre
$A/\mathfrak J^n$) de la aplicación canónica $u_n$; es decir, es
la aplicación canónica de
$\widehat A=\varprojlim_k A/\mathfrak J^k$ sobre $A/\mathfrak J^n$;
$u\circ v$ es, por tanto, límite proyectivo de esta sucesión de aplicaciones,
es decir, por definición la aplicación identidad de $\widehat A$; esto
demuestra la proposición.

\begin{corollary}[7.2.5]
\label{0.7.2.5-fr}
Bajo las hipótesis de \textup{(7.2.3)}, las condiciones siguientes son equivalentes:
\begin{enumerate}
  \item[{\upshape a)}] El homomorfismo $u$ es continuo;

\oldpage[0\textsubscript{I}]{64}

  \item[{\upshape b)}] El homomorfismo $v$ es bicontinuo;
  \item[{\upshape c)}] $A$ es un anillo $\mathfrak J$-ádico.
\end{enumerate}
\end{corollary}

\begin{corollary}[7.2.6]
\label{0.7.2.6-fr}
Sean $A$ un anillo admisible, $\mathfrak J$ un ideal de definición de $A$.
Para que $A$ sea noetheriano, es necesario y suficiente que $A/\mathfrak J$ sea
noetheriano y que $\mathfrak J/\mathfrak J^2$ sea un
$(A/\mathfrak J)$-módulo de tipo finito.
\end{corollary}

Estas condiciones son evidentemente necesarias. Recíprocamente, supongamos que se
verifican; como en virtud de (7.2.4) $A$ es completo para la topología
$\mathfrak J$-preádica, para que sea noetheriano, es necesario y suficiente que
el anillo graduado asociado $\operatorname{grad}(A)$ (para la filtración de los
$\mathfrak J^n$) lo sea ([1], p.~18-07, teorema~4). Ahora bien, sean
$a_1,\ldots,a_n$ unos elementos de $\mathfrak J$ cuyas clases mod.
$\mathfrak J^2$ son generadores de $\mathfrak J/\mathfrak J^2$ como
$A/\mathfrak J$-módulo. Es inmediato por inducción que las clases mod.
$\mathfrak J^{m+1}$ de los monomios de grado total $m$ en los $a_i$
$(1\leq i\leq n)$ forman un sistema de generadores del
$(A/\mathfrak J)$-módulo $\mathfrak J^m/\mathfrak J^{m+1}$. Se concluye que
$\operatorname{grad}(A)$ es un anillo isomorfo a un cociente de
$(A/\mathfrak J)[T_1,\ldots,T_n]$ ($T_i$ indeterminadas), lo que termina la
demostración.

\begin{proposition}[7.2.7]
\label{0.7.2.7-fr}
Sea $(A_i,u_{ij})$ un sistema proyectivo $(i\in\mathbf N)$ de anillos discretos,
y para todo entero $i$, sea $\mathfrak J_i$ el núcleo en $A_i$ del
homomorfismo $u_{0i}:A_i\to A_0$. Se supone que:
\begin{enumerate}
  \item[{\upshape a)}] Para $i\leq j$, $u_{ij}$ es sobreyectivo y su núcleo es
    $\mathfrak J_j^{i+1}$ \textup{(por tanto $A_i$ es isomorfo a
    $A_j/\mathfrak J_j^{i+1}$)}.
  \item[{\upshape b)}] $\mathfrak J_1/\mathfrak J_1^2$ ($=\mathfrak J_1$) es un
    módulo de tipo finito sobre $A_0=A_1/\mathfrak J_1$.
\end{enumerate}
Sea $A=\varprojlim_i A_i$, y para todo entero $n\geq0$, sean $u_n$
el homomorfismo canónico $A\to A_n$, $\mathfrak J^{(n+1)}\subset A$ su
núcleo. En estas condiciones:
\begin{enumerate}
  \item[{\upshape(i)}] $A$ es un anillo ádico, que tiene como ideal de definición
    $\mathfrak J=\mathfrak J^{(1)}$.
  \item[{\upshape(ii)}] Se tiene $\mathfrak J^{(n)}=\mathfrak J^n$ para todo $n\geq1$.
  \item[{\upshape(iii)}] $\mathfrak J/\mathfrak J^2$ es isomorfo a
    $\mathfrak J_1=\mathfrak J_1/\mathfrak J_1^2$, y es, por consiguiente, un módulo
    de tipo finito sobre $A_0=A/\mathfrak J$.
\end{enumerate}
\end{proposition}

La hipótesis de sobreyectividad de los $u_{ij}$ implica que $u_n$ es sobreyectivo;
además, la hipótesis a) implica que $\mathfrak J_j^{j+1}=0$, por tanto $A$ es
un anillo admisible (7.2.2); por definición, los $\mathfrak J^{(n)}$ forman
un sistema fundamental de vecindades de $0$ en $A$, por tanto (ii) implica (i).
Además, se tiene $\mathfrak J=\varprojlim_i\mathfrak J_i$ y las aplicaciones
$\mathfrak J\to\mathfrak J_i$ son sobreyectivas, por tanto (ii) implica (iii), y
se reduce a probar (ii). Por definición, $\mathfrak J^{(n)}$ está formado
por los elementos $(x_k)_{k\geq0}$ de $A$ tales que $x_k=0$ para $k<n$, por tanto
$\mathfrak J^{(n)}\mathfrak J^{(m)}\subset\mathfrak J^{(n+m)}$, es decir,
los $\mathfrak J^{(n)}$ constituyen una \emph{filtración} de $A$. Por otra
parte, $\mathfrak J^{(n)}/\mathfrak J^{(n+1)}$ es isomorfo a la proyección de
$\mathfrak J^{(n)}$ sobre $A_n$; como
$\mathfrak J^{(n)}=\varprojlim_{i>n}\mathfrak J_i^n$, esta proyección no es
otra que $\mathfrak J_n^n$, que es un módulo sobre
$A_0=A_n/\mathfrak J_n$. Sean entonces $a_j=(a_{jk})_{k\geq0}$ unos $r$ elementos de
$\mathfrak J=\mathfrak J^{(1)}$ tales que $a_{11},\ldots,a_{r1}$ forman un
sistema de generadores de $\mathfrak J_1$ sobre $A_0$; vamos a ver que
el conjunto $S_n$ de los monomios de grado total $n$ en los $a_j$ genera el ideal
$\mathfrak J^{(n)}$ de $A$. Como $\mathfrak J_i^{i+1}=0$, es claro, en primer
lugar, que $S_n\subset\mathfrak J^{(n)}$; puesto que $A$ es completo para la
filtración $(\mathfrak J^{(m)})$, basta probar que el conjunto
$\overline S_n$ de las clases mod. $\mathfrak J^{(n+1)}$ de los elementos de $S_n$
genera el módulo graduado $\operatorname{grad}(\mathfrak J^{(n)})$ sobre
el anillo graduado $\operatorname{grad}(A)$ para la filtración anterior
([1], p.~18-06, lema); en virtud de la definición de la multiplicación en
$\operatorname{grad}(A)$,

\oldpage[0\textsubscript{I}]{65}

bastará probar que para todo $m$, $\overline S_m$ es un sistema de
generadores del $A_0$-módulo
$\mathfrak J^{(m)}/\mathfrak J^{(m+1)}$, o, equivalentemente, que $\mathfrak J_m^m$ está
generado por los monomios de grado $m$ en los $a_{jm}$ ($1\leq j\leq r$).
Para ello, queda por mostrar que $\mathfrak J_m$ está (como
$A_m$-módulo) generado por los monomios de grado $\leq m$ con respecto a los
$a_{jm}$; siendo la proposición evidente por definición para $m=1$,
razonemos por inducción sobre $m$, y sea $\mathfrak J'_m$ el sub-$A_m$-módulo
de $\mathfrak J_m$ generado por estos monomios. La relación
$\mathfrak J_{m-1}=\mathfrak J_m/\mathfrak J_m^m$ y la hipótesis de inducción
prueban que $\mathfrak J_m=\mathfrak J'_m+\mathfrak J_m^m$, de donde, puesto que
$\mathfrak J_m^{m+1}=0$, se obtiene
$\mathfrak J_m^m={\mathfrak J'_m}^{\,m}$, y finalmente
$\mathfrak J_m=\mathfrak J'_m$.

\begin{corollary}[7.2.8]
\label{0.7.2.8-fr}
Bajo las condiciones de \textup{(7.2.7)}, para que $A$ sea noetheriano, es necesario y
suficiente que $A_0$ lo sea.
\end{corollary}

Esto resulta inmediatamente de (7.2.6).

\begin{proposition}[7.2.9]
\label{0.7.2.9-fr}
Supongamos verificadas las hipótesis de \textup{(7.2.7)}: para todo entero $i$, sea
$M_i$ un $A_i$-módulo, y, para $i\leq j$, sea
$v_{ij}:M_j\to M_i$ un $u_{ij}$-homomorfismo, tales que $(M_i,v_{ij})$ es un
sistema proyectivo. Supongamos además que $M_0$ es un $A_0$-módulo de tipo
finito, que $v_{ij}$ es sobreyectivo y que su núcleo es
$\mathfrak J_j^{i+1}M_j$. Entonces $M=\varprojlim M_i$ es un $A$-módulo de tipo
finito, y el núcleo del $u_n$-homomorfismo sobreyectivo $v_n:M\to M_n$ es
$\mathfrak J^{n+1}M$ \textup{(de modo que $M_n$ se identifica con
$M/\mathfrak J^{n+1}M=M\otimes_A(A/\mathfrak J^{n+1})$)}.
\end{proposition}

Sean $z_h=(z_{hk})_{k\geq0}$ un sistema de $s$ elementos de $M$ tales que los
$z_{h0}$ ($1\leq h\leq s$) forman un sistema de generadores de $M_0$; se va a
mostrar que los $z_h$ generan el $A$-módulo $M$. El $A$-módulo $M$ es
separado y completo para la filtración formada por los $M^{(n)}$, donde $M^{(n)}$
es el conjunto de los $y=(y_k)_{k\geq0}$ de $M$ tales que $y_k=0$ para $k<n$;
es claro que se tiene $\mathfrak J^{(n)}M\subset M^{(n)}$ y que
$M^{(n)}/M^{(n+1)}=\mathfrak J_n^nM_n$. Se reduce, pues, a mostrar que las
clases de los $z_h$ mod. $M^{(1)}$ generan el módulo graduado
$\operatorname{grad}(M)$ (para la filtración anterior) sobre el anillo graduado
$\operatorname{grad}(A)$ ([1], p.~18-06, lema); para ello, se comprueba
fácilmente que basta de nuevo probar que los $z_{hn}$ ($1\leq h\leq s$)
generan el $A_n$-módulo $M_n$. Se razona de nuevo por inducción sobre
$n$, siendo la proposición evidente por definición para $n=0$; la relación
$M_{n-1}=M_n/\mathfrak J_n^nM_n$ y la hipótesis de inducción muestran que si
$M'_n$ es el submódulo de $M_n$ generado por los $z_{hn}$, se tiene
$M_n=M'_n+\mathfrak J_n^nM_n$, y como $\mathfrak J_n$ es nilpotente, esto
implica $M_n=M'_n$. El mismo razonamiento de paso a los módulos graduados
asociados muestra que la aplicación canónica de $\mathfrak J^{(n)}M$ en
$M^{(n)}$ es sobreyectiva (por tanto biyectiva), es decir, que
$\mathfrak J^{(n)}M=\mathfrak J^nM$ es el núcleo de $M\to M_{n-1}$.

\begin{corollary}[7.2.10]
\label{0.7.2.10-fr}
Sea $(N_i,w_{ij})$ un segundo sistema proyectivo de $A_i$-módulos que verifica las
condiciones de \textup{(7.2.9)}, y sea $N=\varprojlim N_i$. Hay una correspondencia
biunívoca entre los sistemas proyectivos $(h_i)$ de $A_i$-homomorfismos
$h_i:M_i\to N_i$ y los homomorfismos de $A$-módulos $h:M\to N$ (que son
necesariamente continuos para las topologías $\mathfrak J$-ádicas).
\end{corollary}

Es claro que si $h:M\to N$ es un $A$-homomorfismo, se tiene
$h(\mathfrak J^nM)\subset\mathfrak J^nN$, de donde la continuidad de $h$; por
paso a los cocientes, corresponde, pues, a $h$ un sistema proyectivo de
$A_i$-homomorfismos $h_i:M_i\to N_i$, del que $h$ es el límite proyectivo,
de donde el corolario.

\begin{remark}[7.2.11]
\label{0.7.2.11-fr}
Sea $A$ un anillo ádico que tiene un ideal de definición $\mathfrak J$ tal que
$\mathfrak J/\mathfrak J^2$ es un $(A/\mathfrak J)$-módulo de tipo finito;
es claro que los $A_i=A/\mathfrak J^{i+1}$ verifican
\oldpage[0\textsubscript{I}]{66}
las condiciones de (7.2.7); como $A$ es el límite proyectivo de los $A_i$, se
ve que la prop. (7.2.7) da la descripción de \emph{todos} los anillos
ádicos del tipo considerado (y en particular de todos los anillos
\emph{ádicos noetherianos}).
\end{remark}

\begin{example}[7.2.12]
\label{0.7.2.12-fr}
Sean $B$ un anillo, $\mathfrak J$ un ideal de $B$ tal que
$\mathfrak J/\mathfrak J^2$ es un módulo de tipo finito sobre
$B/\mathfrak J$ (o sobre $B$, lo que equivale a lo mismo); pongamos
$A=\varprojlim_n B/\mathfrak J^{n+1}$; $A$ es el separado completado de $B$
provisto de la topología $\mathfrak J$-preádica. Si
$A_n=B/\mathfrak J^{n+1}$, es inmediato que los $A_n$ verifican las
condiciones de (7.2.7); por tanto $A$ es un anillo ádico y si
$\overline{\mathfrak J}$ es la adherencia en $A$ de la imagen canónica de
$\mathfrak J$, $\overline{\mathfrak J}$ es un ideal de definición de $A$,
$\overline{\mathfrak J}^n$ es la adherencia de la imagen canónica de
$\mathfrak J^n$, $A/\overline{\mathfrak J}^n$ se identifica con
$B/\mathfrak J^n$ y
$\overline{\mathfrak J}/\overline{\mathfrak J}^2$ es isomorfo a
$\mathfrak J/\mathfrak J^2$ como
$(A/\overline{\mathfrak J})$-módulo. Del mismo modo, si $N$ es tal que
$N/\mathfrak JN$ es un $B$-módulo de tipo finito, y si se pone
$M_i=N/\mathfrak J^{i+1}N$, $M=\varprojlim M_i$ es un $A$-módulo de tipo
finito, isomorfo al separado completado de $N$ para la topología
$\mathfrak J$-preádica, $\overline{\mathfrak J}^nM$ se identifica con
la adherencia de la imagen canónica de $\mathfrak J^nN$, y
$M/\overline{\mathfrak J}^nM$ con $N/\mathfrak J^nN$.
\end{example}

